% !TEX spellcheck = English

\documentclass[11pt, a4paper]{amsart}

\usepackage{amsmath,amsthm,amssymb}

\newcommand{\Z}{\mathbb{Z}}
\newcommand{\Zp}{\mathbb{Z}^{+}}

\theoremstyle{plain}
\newtheorem{thm}{Theorem}
% \newtheorem{cor}[thm]{Corollary}
\newtheorem*{roth}{Roth's theorem}

\theoremstyle{remark}
\newtheorem*{rem}{Remark}

%% PROVISIONAL, REMOVE:
%\usepackage{color, soul}
%\usepackage[notcite,notref]{showkeys}

\title{The continued fraction expansion of certain Pierce series}
\author{Juan Luis Varona}
%\date{} % Activate to display a given date or no date

\address{Departamento de Matem\'aticas y Computaci\'on\\
    Universidad de La Rioja\\
    26006 Logro\~no, Spain}
\email{jvarona@unirioja.es}
\urladdr{http://www.unirioja.es/cu/jvarona/}

\subjclass[2010]{Primary: 11J70; Secondary: 11B37, 11J81}

\keywords{Continued fractions, Pierce series, Transcendental numbers}

\thanks{The research of the author is supported by grant MTM2015-65888-C4-4-P of MINECO/FEDER}

\renewcommand{\datename}{\textbf{This paper has been published in}:}
\date{\textit{J.\@ Number Theory} \textbf{180} (2017), 573--578.}


%----------------
\begin{document}
%----------------

\begin{abstract}
A Pierce series is an alternating sum of the reciprocals of an increasing sequence of positive integers, each one divisible by the previous one. In this short note we study some Pierce series that arise from certain continued fractions. Moreover, we show that the sums of these Pierce series are transcendental numbers. 
\end{abstract}

%----------------
\maketitle
%----------------

%----------------
\section{Introduction and main results}
\label{sec:intro}
%----------------

Given an eventually increasing sequence of positive integers $(x_n)$ such that $x_n \mid x_{n+1}$ for all $n$, the sum of the reciprocals is the Engel series
\begin{equation}
\label{eq:Engel}
  \sum_{j=1}^{\infty} \frac{1}{x_j}
  = \sum_{j=1}^{\infty} \frac{1}{y_1y_2 \cdots y_j},
\end{equation}
where $y_1 = x_1$ and $y_{n+1} = x_{n+1}/x_n$ for $n \ge 1$,
and the alternating sum of the reciprocals is the Pierce series
\begin{equation}
\label{eq:Pierce}
  \sum_{j=1}^{\infty} \frac{(-1)^{j+1}}{x_j}
  = \sum_{j=1}^{\infty} \frac{(-1)^{j+1}}{y_1y_2 \cdots y_j}.
\end{equation}

In the recent papers \cite{Hone1, Hone2}, A.N.W. Hone gives some families of sequences $(x_n)$ that are generated by certain nonlinear recurrences of second order, and such that some continued fraction expansion defined from the $(x_n)$ coincides with the Engel series~\eqref{eq:Engel}. Additionally, he proves, under some circumstances, that~\eqref{eq:Engel} is a transcendental number. 
We give a similar result for Pierce series, although the structure of the corresponding continued fractions is different.

To generate the sequences $(x_n)$, let us take the initial values $x_0 = x_1 = 1$ and define $(x_n)$ by the recurrence relation
\begin{equation}
\label{eq:rec}
  x_{n+2}x_n = (1+\alpha_{n+1}x_{n+1})x_{n+1}^2, \quad n \ge 0,
\end{equation}
with $\alpha_{n+1}$ any positive integer.
If we take
\begin{equation}
\label{eq:xyz}
  y_n = \frac{x_{n+1}}{x_n},
  \qquad
  z_n = \frac{y_{n+1}}{y_n} = \frac{x_{n+2}x_n}{x_{n+1}^2} = 1+\alpha_{n+1}x_{n+1},
\end{equation}
it is clear that $y_0=1$ and $x_2=y_1=z_0=1+\alpha_1$; moreover, by induction we can see that $x_n,y_n,z_n \in \Zp = \{1,2,3,\dots\}$ (use~\eqref{eq:xyz} to prove that, if $x_{n+2},y_{n+1},z_n \in \Zp$, then $z_{n+1},y_{n+2},x_{n+3} \in \Zp$). 

We use the standard notation \cite{Cas, Hua, Khin, NZM, Var} for continued fractions, namely 
\[
  [a_0;a_1,a_2,\dots,a_n]
  = a_0 + \cfrac{1}{ a_1 + \cfrac{1}{ a_2 + \cdots \cfrac{1}{a_n} } }
  = \frac{p_n}{q_n},
\]
with $a_0 \in \Z$ and $a_j \in \Zp$ for $j \ge 1$ (we also assume $q_n>0$), and 
\[
  [a_0;a_1,a_2,\dots] = \lim_{n\to\infty} [a_0;a_1,a_2,\dots,a_n],
\]
a limit that always exists.

The convergents $p_n/q_n$ of a continued fraction satisfy the recurrence relation
\begin{equation}
\label{eq:fcpkqk}
\begin{aligned}
  p_{-1} &= 1, & p_0 &= a_0, & p_1 &= a_0a_1+1, & p_k &= a_kp_{k-1}+p_{k-2},\ k \ge 1,\\
  q_{-1} &= 0, & q_0 &= 1, & q_1 &= a_1, & q_k &= a_kq_{k-1}+q_{k-2},\ k \ge 1.
\end{aligned}
\end{equation}
Among the large number of well-known properties of continued fractions, let us mention that
they satisfy 
\[
  p_k q_{k-1} - p_{k-1} q_k = (-1)^{k-1}, \qquad k \ge 1.
\]
Using this property (written as $p_k/q_k - p_{k-1}/q_{k-1} = (-1)^{k-1}/(q_{k-1}q_k)$) and induction, we obtain
\begin{equation}
\label{eq:fcaltqk}
  [a_0;a_1,a_2,\dots,a_n]
  = a_0 + \frac{1}{q_1} - \frac{1}{q_1q_2} + \frac{1}{q_2q_3}
  + \cdots + \frac{(-1)^{n-1}}{q_{n-1}q_n}.
\end{equation}
This provides, of course, an alternating series, but it is not a Peirce
series because, in general, $q_j \nmid q_{j+2}$
(actually, this is what happens in~\cite{DaSha}, another paper that
studies continued fractions for some alternating series whose
sum is a transcendental number).
Anyway, the existence of this alternating series closely related
to the basic theory of continued fractions is a strong motivation for searching for 
series of type~\eqref{eq:Pierce} (Pierce series), perhaps more than
series of type~\eqref{eq:Engel} (Engel series).
Let us also remember that alternating series (with the general term
decreasing to zero) have a great interest due to their very good approximation 
properties: according to Leibniz's theorem, any partial sum approximates the
sum of the series with an error bounded by the first omitted term.

The proof of the following theorem can be found in Section~\ref{sec:proof}.

\begin{thm}
\label{th:main}
Let $(x_n)$ be a sequence generated from the initial values $x_0 = x_1 = 1$ by the recurrence~\eqref{eq:rec}. The partial alternating sums of their reciprocals have the continued fraction expansions
\begin{equation}
\label{eq:SN}
  S_N := \sum_{j=1}^N \frac{(-1)^{j+1}}{x_j} 
  = [a_0;a_1,a_2,\dots,a_{3N-4}] 
\end{equation}
for all $N \ge 2$, where $a_0=0$, $a_1=1$, $a_2=x_2-1$, $a_3 = \alpha_2x_2$, and
\begin{equation}
\label{eq:an}
  a_{3n+1} = 1, \quad
  a_{3n+2} = x_{n+1}-1, \quad
  a_{3n+3} = \frac{\alpha_{n+2}x_{n+2}}{x_{n+1}} - 1 \in \Zp,
  \quad n \ge 1.
\end{equation}
\end{thm}

It is interesting to note that this result is somewhat different to what happens in the case of Engel series~\eqref{eq:Engel}, where we have a relation of the type
\[
  \sum_{j=1}^N \frac{1}{x_j} 
  = [b_0;b_1,b_2,\dots,b_{2N-2}],
\]
and where, instead of \eqref{eq:an}, only two expressions $b_{2n} = x_n$ and $b_{2n+1} = \alpha_{n+1}x_{n+1}/x_n$ appear. The details can be found in~\cite{Hone1}.

Finally, taking $N \to \infty$ in Theorem~\ref{th:main} and using Roth's theorem on rational approximation and transcendence, we have the following result whose proof can be seen in Section~\ref{sec:transc}:

\begin{thm}
\label{th:transc}
The infinite sum 
\begin{equation}
\label{eq:Sinfty}
  \sum_{j=1}^\infty \frac{(-1)^{j+1}}{x_j} 
  = [a_0;a_1,a_2,\dots,a_{n},\dots]
\end{equation}
(where the coefficients of the continued fraction are as in~\eqref{eq:an})
is a transcendental number.
\end{thm}

%----------------
\section{Proof of Theorem~\ref{th:main}}
\label{sec:proof}
%----------------

Let us take the sequence $(a_n)$ as in~\eqref{eq:an}, whose terms belong to $\Zp$ for $n>0$
because we can write $a_{3n+3} = \alpha_{n+2}x_{n+2}/x_{n+1}-1 = y_{n+1}-1$.
% Indeed, it is enough to notice that we can write $a_{3n+3} = \alpha_{n+2}x_{n+2}/x_{n+1}-1 = y_{n+1}-1$.

Now, let us denote $\frac{p_n}{q_n} = [a_0;a_1,a_2,\dots,a_n]$. 
From $a_0 = p_0/q_0$, $a_0+1/a_1 = p_1/q_1$, $a_0+1/(a_1+1/a_2) = p_2/q_2$, and taking into account that $a_0=0$, $a_1=1$, $a_2=x_2-1$, we easily obtain $p_0=0$, $q_0=1$, $p_1=1$, $q_1=a_1=1$, $p_2=a_2=x_2-1$ and $q_2=a_2+1=x_2$.
Let us prove now by induction that the denominators $q_n$ satisfy the recurrence relation
\begin{equation}
\label{eq:qn}
  q_{3n-1} = x_{n+1}, \quad
  q_{3n} = \frac{x_{n+2}}{x_{n+1}}-x_{n+1}+1, \quad
  q_{3n+1} = \frac{x_{n+2}}{x_{n+1}}+1, % \quad
%  q_{3n+2} = x_{n+2}, \quad
%  q_{3n+3} = \frac{x_{n+3}}{x_{n+2}}-x_{n+2}+1 \in \Zp,
  \quad n \ge 1.
\end{equation}

For $n=1$ we have already seen that $q_2=x_2$; 
by using \eqref{eq:fcpkqk} and $(1+\alpha_2x_2)x_2 = x_1x_3/x_2$ (with $x_1=1$), we get
$q_3 = a_3q_2 + q_1 = (-1+1+\alpha_2x_2)x_2+1 = x_3/x_2-x_2+1$; 
and, finally, 
$q_4 = a_4q_3 + q_2 = 1(x_3/x_2-x_2+1)+x_2 = x_3/x_2+1$.
Let us now check the induction step.
Assuming \eqref{eq:qn} for $n$, for $n+1$ we have
\begin{align*}
q_{3n+2} &= a_{3n+2} q_{3n+1} + q_{3n} \\
&= (x_{n+1}-1) \Big( \frac{x_{n+2}}{x_{n+1}}+1 \Big)
+ \frac{x_{n+2}}{x_{n+1}}-x_{n+1}+1
= x_{n+2},\\
q_{3n+3} &= a_{3n+3} q_{3n+2} + q_{3n+1} \\
&= \Big( \frac{-1+1+\alpha_{n+2}x_{n+2}}{x_{n+1}} - 1 \Big) x_{n+2}
+ \frac{x_{n+2}}{x_{n+1}}+1
= \frac{x_{n+3}}{x_{n+2}}-x_{n+2}+1,\\
q_{3n+4} &= a_{3n+4} q_{3n+3} + q_{3n+2} \\
&= 1 \Big( \frac{x_{n+3}}{x_{n+2}} - x_{n+2} + 1 \Big) 
+ x_{n+2} = \frac{x_{n+3}}{x_{n+2}}+1.
\end{align*}

Now, let us prove~\eqref{eq:SN} by induction. 
For $N=2$, \eqref{eq:SN} becomes
\[
  S_2 = \frac{1}{x_1} - \frac{1}{x_2} 
  = [a_0;a_1,a_2],
\]
which follows easily from the equalities $a_0=0$, $a_1=1$, $a_2=x_2-1$.
For the induction step, let us use~\eqref{eq:fcaltqk}, which allows us
to write
\[
  [a_0;a_1,a_2,\dots,a_{3N-4}]
  = a_0 + \frac{1}{q_1} - \frac{1}{q_1q_2} + \frac{1}{q_2q_3}
  + \cdots + \frac{(-1)^{3N-5}}{q_{3N-5}q_{3N-4}}
\]
and the corresponding expression for $[a_0;a_1,a_2,\dots,a_{3N-1}]$.
Then, we only need to check that
\[
  \frac{1}{x_{N+1}} = \frac{1}{q_{3N-4}q_{3N-3}}
  - \frac{1}{q_{3N-3}q_{3N-2}} 
  + \frac{1}{q_{3N-2}q_{3N-1}}.
\]
According to~\eqref{eq:qn}, we have
\begin{align*}
  & \frac{1}{q_{3N-4}q_{3N-3}}
  - \frac{1}{q_{3N-3}q_{3N-2}} 
  + \frac{1}{q_{3N-2}q_{3N-1}} \\
  &\qquad = 
  \frac{1}{x_{N+1}-x_N^2+x_N}
  - \frac{x_N^2}{(x_{N+1}-x_N^2+x_N)(x_{N+1}+x_N)} \\*
  &\qquad\qquad\qquad+ \frac{x_N}{(x_{N+1}+x_N)x_{N+1}} \\
  &\qquad =
  \frac{x_{N+1} (x_{N+1}+x_N) - x_N^2x_{N+1} + x_N (x_{N+1}-x_N^2+x_N)}
    {(x_{N+1}-x_N^2+x_N)(x_{N+1}+x_N)x_{N+1}} \\
  &\qquad =
  \frac{x_{N+1} (x_{N+1}-x_N^2+x_N) + x_N (x_{N+1}-x_N^2+x_N)}
    {(x_{N+1}-x_N^2+x_N)(x_{N+1}+x_N)x_{N+1}} 
  = \frac{1}{x_{N+1}},
\end{align*}
and the proof is complete.

%----------------
\section{Proof of Theorem~\ref{th:transc}}
\label{sec:transc}
%----------------

The transcendence of~\eqref{eq:Sinfty} is due to the large rate of growth of the sequence~$(x_n)$. The analysis of the increasing of $(x_n)$ is already done in \cite{Hone1, Hone2}, but we repeat it here for the sake of completeness. By definition, $\alpha_n \ge 1$, so that
\begin{equation}
\label{eq:n+1n}
  x_{n+1} = (1+\alpha_nx_n)x_n^2/x_{n-1} > x_n^3/x_{n-1} \ge x_n^2,
  \qquad n \ge 1.
\end{equation}
This rate of growth is not enough for our purposes, but this result can be easily improved. Let us write~\eqref{eq:n+1n} as $x_{n-1} \le x_n^{1/2}$, and use it in $x_{n+1} > x_n^3/x_{n-1}$. 
In this way, we obtain
\begin{equation}
\label{eq:2+1/2}
  x_{n+1} > x_n^3/x_{n-1} \ge x_n^{2+1/2}, 
  \qquad n \ge 1.
\end{equation}

\begin{roth}[1955, \cite{Roth, Cas}]
Let $\xi$ be an irrational algebraic number. Then, for any arbitrary fixed $\varepsilon > 0$, there are only finitely many rational approximations~$p/q$ (we assume $q>0$) for which
\[
  \left| \xi - \frac{p}{q} \right| < \frac{1}{q^{2+\varepsilon}}.
\]
\end{roth}

Now, Theorem~\ref{th:transc} follows immediately from Leibnitz's error estimate for the truncation of alternating series, \eqref{eq:2+1/2}, and Roth's theorem:
\[
  \left| \sum_{j=1}^\infty \frac{(-1)^{j+1}}{x_j} - \frac{p_{3n-1}}{q_{3n-1}} \right| 
  = \left| \sum_{j=n+2}^{\infty} \frac{(-1)^{j+1}}{x_j} \right|
  < \frac{1}{x_{n+2}} < \frac{1}{x_{n+1}^{2+1/2}} = \frac{1}{q_{3n-1}^{2+1/2}}.
\]

%----------------
\begin{thebibliography}{0}
%----------------

\bibitem{Cas}
J. W. S. Cassels, 
\textit{An Introduction to Diophantine Approximation}, 
Cambridge University Press, 1957.

\bibitem{DaSha}
J. L. Davison and J. O. Shallit, 
Continued fractions for some alternating series, 
\textit{Monatsh. Math.} \textbf{111} (1991), 119--126.

\bibitem{Hone1}
A. N. W. Hone, 
Curious continued fractions, nonlinear recurrences and transcendental numbers, 
\textit{J. Integer Seq.} \textbf{18} (2015), no.~8, article 15.8.4, 10~pp.

\bibitem{Hone2}
A. N. W. Hone, 
On the continued fraction expansion of certain Engel series, 
\textit{J. Number Theory} \textbf{164} (2016), 269--281.

\bibitem{Hua}
Hua Loo Keng, 
\textit{Introduction to Number Theory}, 
Springer-Verlag, 1982.

\bibitem{Khin}
A. Ya. Khinchin, 
\textit{Continued Fractions}, 
University of Chicago Press, 1964.
Reprint: Dover, 1997.

\bibitem{NZM}
I. Niven, H. S. Zuckerman and H. L. Montgomery, 
\textit{An introduction to the Theory of Numbers}, 5th ed., 
Wiley, 1991.

\bibitem{Roth}
K. F. Roth, 
Rational approximation of algebraic numbers, 
\textit{Mathematika} \textbf{2} (1955), 1--20; 
Corrigendum, \textit{Mathematika} \textbf{2} (1955), 168.

\bibitem{Var}
J. L. Varona, 
\textit{Recorridos por la Teor\'ia de N\'umeros} (Spanish),
Electolibris \& Real Sociedad Matem\'atica Espa\~nola, Murcia, 2014.

%----------------
\end{thebibliography}
%----------------

%----------------
\end{document}
%----------------
